Online web agents often augment a base actor with memory, workflow, or skill modules.
These modules can improve performance, but they also consume test-time tokens, a cost rarely reported alongside the actor's inference cost.
We study online augmentation, where this overhead is paid on every task, and re-evaluate its benefits under a fixed total inference budget.
We compare AWM, ASI, and ReasoningBank with a token-matched vanilla baseline that uses the same budget for additional actor steps.
Across four WebArena domains and three models, Gemini~3~Flash, GPT-5.4-mini, and Qwen~3.6-27B, the vanilla baseline matches or surpasses all three augmentation methods in aggregate success rate while often using fewer total tokens.
We observe a similar trend on WorkArena-L1 with Qwen~3.6-27B, indicating that the effect extends to enterprise knowledge-work tasks.
Our results suggest that skills and workflow memory can be useful in specific domains, but their apparent gains often vanish against a budget-matched actor.
We further show that run-to-run variance materially affects outcomes and should be reported as a core evaluation criterion for online web agents.